%----------------------------------------------------------------------------------------
% Lettre de motivation — modèle find-job (une page, pdflatex, paquets standard)
% Pour utiliser votre propre modèle : placez lettre.tex à la racine de l'espace.
%----------------------------------------------------------------------------------------
\documentclass[11pt, a4paper]{article}

\usepackage[T1]{fontenc}
\usepackage[utf8]{inputenc}
\usepackage[french]{babel}
\usepackage[top=2cm, bottom=2cm, left=2.3cm, right=2.3cm]{geometry}
\usepackage{lmodern}
\usepackage{microtype}
\usepackage{xcolor}
\usepackage{hyperref}

\definecolor{bleuencre}{RGB}{22,38,96}
\definecolor{darkgray}{RGB}{50,50,50}
\definecolor{midgray}{RGB}{105,105,105}
\hypersetup{colorlinks=true, urlcolor=bleuencre}
\pagestyle{empty}
\color{darkgray}
\setlength{\parindent}{0pt}
\setlength{\parskip}{0.8em}
\linespread{1.08}

%----------------------------------------------------------------------------------------
% INFORMATIONS — renseignées par find-job prepare, à compléter par l'agent
%----------------------------------------------------------------------------------------
\newcommand{\candidat}{Camille Martin}
\newcommand{\coordonnees}{camille.martin@example.com \quad 06 12 34 56 78 \quad Lyon}
\newcommand{\entreprise}{Société Générale}
\newcommand{\destinataire}{Équipe recrutement}
\newcommand{\poste}{Data Engineer Senior}
\newcommand{\datelettre}{Le 24 octobre 2026}

\begin{document}

{\Large\bfseries\color{bleuencre}\candidat}\par\vspace{-0.6em}
{\small\color{midgray}\coordonnees}

\vspace{1.2em}
\begin{flushright}
  \textbf{\entreprise}\\
  \destinataire\\[0.8em]
  \datelettre
\end{flushright}

\textbf{Objet :} candidature -- \poste

Madame, Monsieur,

Les équipes data de Société Générale traitent des volumes que peu d'entreprises connaissent, et la fiabilité des pipelines y compte autant que leur vitesse. C'est ce qui m'attire dans votre offre de Data Engineer Senior.

Chez Logistique Rhône, je conçois depuis 2022 les pipelines d'ingestion qui reçoivent chaque jour deux millions d'événements. J'ai migré notre entrepôt vers PostgreSQL et dbt, ce qui a réduit le temps de calcul de 40\,\%, puis j'ai ajouté des tests de qualité et des alertes sur trente tables critiques. Avant cela, à la Banque Saône, j'ai automatisé douze rapports hebdomadaires pour l'équipe marketing : je connais donc aussi les attentes des métiers bancaires.

J'aimerais apporter cette rigueur à vos flux Airflow et découvrir à vos côtés le traitement en temps réel avec Kafka, que je n'ai encore pratiqué que sur des projets personnels.

Je serais heureuse d'en parler avec vous lors d'un échange.

Je vous prie d'agréer, Madame, Monsieur, mes salutations distinguées.

\vspace{0.6em}
\candidat

\end{document}
